\documentclass[12pt, a4paper]{article}
\usepackage{amsmath, amssymb, amsthm}
\usepackage{geometry}
\usepackage{hyperref}
\usepackage{bm}
\usepackage{enumitem}
\usepackage{titlesec}
\usepackage{xcolor}

\geometry{margin=2.5cm}

\title{\textbf{Chapter 3 Summary}\\[0.3em]
\large Kalman Filter and Its Variants}
\author{ESE 650 --- Learning in Robotics}
\date{Spring 2026}

\begin{document}
\maketitle
\tableofcontents
\newpage

%============================================================
\section{Background: Gaussian Random Variables (Section 3.1)}
%============================================================

\subsection{Covariance Matrix}

For a $d$-dimensional random variable $X\in\mathbb{R}^d$:
\begin{equation}
  \Sigma = \mathrm{Cov}(X) = \mathrm{E}[(X-\mathrm{E}[X])(X-\mathrm{E}[X])^\top] \in \mathbb{R}^{d\times d}.
\end{equation}
$\Sigma$ is symmetric positive semi-definite: $\Sigma=U\Lambda U^\top$
with $\Lambda\ge0$. Uncertainty: $\mathrm{tr}(\Sigma)=\sum_i\lambda_i(\Sigma)$.

Useful trace identities: $\mathrm{tr}(AB)=\mathrm{tr}(BA)$; $\mathrm{tr}(A^\top B)=\sum_{ij}B_{ij}A_{ij}$.

\subsection{Gaussian Distribution}

The $d$-dimensional \textbf{Gaussian} $\mathcal{N}(\mu,\Sigma)$ has PDF:
\begin{equation}
  f(x) = \frac{1}{\sqrt{\det(2\pi\Sigma)}}
  \exp\!\left\{-\tfrac{1}{2}(x-\mu)^\top\Sigma^{-1}(x-\mu)\right\}.
\end{equation}

\textbf{Linear transform}: if $X\sim\mathcal{N}(\mu_X,\Sigma_X)$ and
$Y=AX$ (deterministic $A$), then
\begin{equation}
  Y\sim\mathcal{N}(A\mu_X,\; A\Sigma_X A^\top).
  \label{eq:gaussian_linear}
\end{equation}
This result is fundamental to both the KF and EKF.

%============================================================
\section{Linear State Estimation (Section 3.2)}
%============================================================

\subsection{Optimally Combining Two Estimators}

Given two unbiased estimators $\hat{X}_1\sim\mathcal{N}(X,\Sigma_1)$
and $\hat{X}_2\sim\mathcal{N}(X,\Sigma_2)$, combine linearly:
$\hat{X}=K_1\hat{X}_1+K_2\hat{X}_2$.
Unbiasedness requires $K_1+K_2=I$.
Minimising $\mathrm{tr}(\Sigma_{\hat{X}})$ gives:
\begin{equation}
  \hat{X} = \Sigma_2(\Sigma_1+\Sigma_2)^{-1}\hat{X}_1
           + \Sigma_1(\Sigma_1+\Sigma_2)^{-1}\hat{X}_2.
  \label{eq:opt_combine}
\end{equation}

\textbf{Key insight:} combining two estimators always yields a better
(lower variance) estimator; combining never hurts.

\subsection{Incorporating a Gaussian Observation}

Given prior estimate $\hat{X}'$ and linear observation
$Y=CX+\nu$, $\nu\sim\mathcal{N}(0,Q)$, the updated estimator is:
\begin{equation}
  \hat{X} = \hat{X}' + K(Y - C\hat{X}'),
  \label{eq:update}
\end{equation}
where $Y-C\hat{X}'$ is called the \textbf{innovation}. The optimal
gain (\textbf{Kalman gain}) that minimises $\mathrm{tr}(\Sigma_{\hat{X}})$:
\begin{equation}
  \boxed{K = \Sigma_{\hat{X}'} C^\top (C\Sigma_{\hat{X}'} C^\top + Q)^{-1}.}
  \label{eq:Kgain}
\end{equation}

The updated covariance:
\begin{align}
  \Sigma_{\hat{X}} &= (I-KC)\Sigma_{\hat{X}'}(I-KC)^\top + KQK^\top
  \quad\text{(Joseph's form, numerically stable)} \notag\\
  &= \left(\Sigma_{\hat{X}'}^{-1} + C^\top Q^{-1}C\right)^{-1}.
  \label{eq:Sigma_update}
\end{align}

%============================================================
\section{Dynamical Systems (Section 3.3)}
%============================================================

\subsection{Linear Time-Invariant (LTI) Systems}

Continuous-time state-space model:
\begin{equation}
  \dot{x}(t) = Ax(t) + Bu(t), \qquad y(t) = Cx(t) + Du(t).
  \label{eq:LTI_ct}
\end{equation}

Discrete-time version:
\begin{equation}
  x_{k+1} = A_{\Delta t} x_k + B_{\Delta t} u_k, \qquad y_k = Cx_k + Du_k.
  \label{eq:LTI_dt}
\end{equation}

\subsection{Nonlinear Systems}

\begin{equation}
  x_{k+1} = f(x_k,u_k) + \epsilon_k, \qquad y_k = g(x_k,u_k) + \nu_k.
  \label{eq:nonlinear}
\end{equation}

\subsection{Markov Decision Processes (MDPs)}

An MDP models systems where the dynamics $f(x_k,u_k)$ is not precisely
known. The state evolves stochastically:
\begin{equation}
  x_{k+1} \sim \mathrm{P}(x_{k+1}\mid x_k, u_k), \qquad
  x_{k+1} = Ax_k + Bu_k + \epsilon_k \quad\text{(linear stochastic)}.
\end{equation}
Noise $\epsilon_k$ is known in distribution, independent across timesteps,
and is not under our control.

%============================================================
\section{Kalman Filter (Section 3.5)}
%============================================================

\subsection{System Model}

Linear stochastic system with linear observations:
\begin{equation}
  x_{k+1} = Ax_k + Bu_k + \epsilon_k, \qquad
  y_k = Cx_k + \nu_k,
  \label{eq:KF_model}
\end{equation}
where $\epsilon_k\sim\mathcal{N}(0,R)$, $\nu_k\sim\mathcal{N}(0,Q)$ are
independent white noise processes.

\textbf{Key insight}: since all operations are linear and Gaussian, the
filtering distribution stays Gaussian at every step:
$\hat{x}_{k|k}\sim\mathcal{N}(\mu_{k|k},\Sigma_{k|k})$.

\subsection{Step 1 --- Propagate Dynamics}

Given estimate $\hat{x}_{k|k}\sim\mathcal{N}(\mu_{k|k},\Sigma_{k|k})$,
predict the state at $k+1$ before the next observation:
\begin{align}
  \mu_{k+1|k} &= A\mu_{k|k} + Bu_k, \label{eq:KF_prop_mu}\\
  \Sigma_{k+1|k} &= A\Sigma_{k|k}A^\top + R. \label{eq:KF_prop_Sigma}
\end{align}

\subsection{Step 2 --- Incorporate Observation}

Upon receiving $y_{k+1}=Cx_{k+1}+\nu_{k+1}$, update using the Kalman
gain:
\begin{align}
  K_{k+1} &= \Sigma_{k+1|k}\,C^\top\!\left(C\Sigma_{k+1|k}C^\top+Q\right)^{-1},
  \label{eq:KF_gain}\\
  \mu_{k+1|k+1} &= \mu_{k+1|k} + K_{k+1}(y_{k+1}-C\mu_{k+1|k}),
  \label{eq:KF_update_mu}\\
  \Sigma_{k+1|k+1} &= (I-K_{k+1}C)\,\Sigma_{k+1|k}.
  \label{eq:KF_update_Sigma}
\end{align}

\subsection{Properties of the KF}

\begin{itemize}
  \item \textbf{Optimality}: among all estimators (linear or nonlinear),
    the KF achieves the minimum mean-squared error for linear Gaussian
    systems. No other method can do better.
  \item \textbf{Best linear filter}: even for non-Gaussian noise, the KF
    is the best \emph{linear} filter.
  \item \textbf{Assumptions}: dynamics noise $\epsilon_k$ and observation
    noise $\nu_k$ are white (uncorrelated in time) and mutually
    independent.
\end{itemize}

%============================================================
\section{Extended Kalman Filter (Section 3.6)}
%============================================================

\subsection{Linearisation of Nonlinear Functions}

For a nonlinear $f:\mathbb{R}^d\to\mathbb{R}^p$ and $x\sim\mathcal{N}(\mu_x,\Sigma_x)$,
linearise around $\mu_x$:
\begin{equation}
  y \approx f(\mu_x) + J(x-\mu_x), \qquad J=\left.\frac{\mathrm{d}f}{\mathrm{d}x}\right|_{x=\mu_x}.
  \label{eq:taylor}
\end{equation}
This gives the first-order approximation:
\begin{equation}
  \mathrm{E}[y]\approx f(\mu_x), \qquad \Sigma_y\approx J\Sigma_x J^\top.
  \label{eq:EKF_propagation}
\end{equation}

\subsection{EKF Algorithm}

For the nonlinear system $x_{k+1}=f(x_k,u_k)+\epsilon_k$,
$y_k=g(x_k)+\nu_k$:

\textbf{Step 1 --- Propagate (linearise dynamics at $\mu_{k|k}$):}
\begin{align}
  A(\mu_{k|k}) &= \left.\frac{\partial f}{\partial x}\right|_{x=\mu_{k|k}}, \\
  \mu_{k+1|k} &= f(\mu_{k|k}, u_k), \\
  \Sigma_{k+1|k} &= A\Sigma_{k|k}A^\top + R.
\end{align}

\textbf{Step 2 --- Update (linearise observations at $\mu_{k+1|k}$):}
\begin{align}
  C(\mu_{k+1|k}) &= \left.\frac{\partial g}{\partial x}\right|_{x=\mu_{k+1|k}}, \\
  K &= \Sigma_{k+1|k}C^\top(C\Sigma_{k+1|k}C^\top+Q)^{-1}, \\
  \mu_{k+1|k+1} &= \mu_{k+1|k} + K\!\left(y_{k+1}-g(\mu_{k+1|k})\right), \\
  \Sigma_{k+1|k+1} &= (I-KC)\Sigma_{k+1|k}.
\end{align}

The EKF replaces the constant matrices $A,C$ in the KF with Jacobians
recomputed at each timestep.

%============================================================
\section{Unscented Kalman Filter (Section 3.7)}
%============================================================

\subsection{Motivation}

The EKF approximates $\mathrm{P}(f(x))$ as Gaussian by linearising $f$
around the mean. This can be inaccurate when $f$ is highly nonlinear.
The \textbf{Unscented Transform} (UT) achieves accuracy to \emph{third
order} vs.\ first order for linearisation, without computing Jacobians.

\subsection{Unscented Transform (Section 3.7.1)}

Given $x\sim\mathcal{N}(\mu,\Sigma)\in\mathbb{R}^n$, choose $2n$
\textbf{sigma points} with equal weights $w^{(i)}=\frac{1}{2n}$:
\begin{equation}
  x^{(i)} = \mu + (\sqrt{n\Sigma})_i^\top, \qquad
  x^{(n+i)} = \mu - (\sqrt{n\Sigma})_i^\top, \quad i=1,\ldots,n.
  \label{eq:sigma_pts}
\end{equation}

Transform each sigma point: $y^{(i)}=f(x^{(i)})$.

Compute the mean and covariance of the transformed distribution:
\begin{align}
  \mu_y &= \sum_{i=1}^{2n} w^{(i)}\,y^{(i)}, \\
  \Sigma_y &= \sum_{i=1}^{2n} w^{(i)}\,(y^{(i)}-\mu_y)(y^{(i)}-\mu_y)^\top.
  \label{eq:UT}
\end{align}

\subsection{UKF Algorithm (Section 3.7.3)}

\textbf{Step 1 --- Propagate using UT:}
\begin{align}
  \mu_{k+1|k} &:= \sum_{i=1}^{2n} w^{(i)}\,f(x^{(i)},u_k), \\
  \Sigma_{k+1|k} &:= R + \sum_{i=1}^{2n} w^{(i)}\,(f(x^{(i)},u_k)-\mu_{k+1|k})(f(x^{(i)},u_k)-\mu_{k+1|k})^\top.
\end{align}

\textbf{Step 2 --- Update using UT:} compute new sigma points for
$\mathcal{N}(\mu_{k+1|k},\Sigma_{k+1|k})$, then:
\begin{align}
  \hat{y} &= \sum_{i}w^{(i)}\,g(x^{(i)}), \\
  \Sigma_{yy} &:= Q + \sum_i w^{(i)}(g(x^{(i)})-\hat{y})(g(x^{(i)})-\hat{y})^\top, \\
  \Sigma_{xy} &:= \sum_i w^{(i)}(x^{(i)}-\mu_{k+1|k})(g(x^{(i)})-\hat{y})^\top.
\end{align}

Kalman gain and update:
\begin{align}
  K^* &= \Sigma_{xy}\Sigma_{yy}^{-1}, \\
  \mu_{k+1|k+1} &= \mu_{k+1|k} + K^*(y_{k+1}-\hat{y}), \\
  \Sigma_{k+1|k+1} &= \Sigma_{k+1|k} - K^*\Sigma_{yy}K^{*\top}.
\end{align}

\subsection{EKF vs.\ UKF}

\begin{itemize}
  \item \textbf{EKF}: requires Jacobians, first-order accuracy,
    simpler to implement.
  \item \textbf{UKF}: no Jacobians, third-order accuracy in mean and
    covariance, more involved implementation.
  \item \textbf{Equivalence}: for linear systems, both reduce to the KF.
  \item \textbf{Limitation}: both still approximate the filtering
    distribution as a \emph{Gaussian}.
\end{itemize}

%============================================================
\section{Particle Filters (Section 3.8)}
%============================================================

\subsection{Motivation}

KF/EKF/UKF all maintain a Gaussian approximation. For highly nonlinear
or non-Gaussian systems, a richer representation is needed. Particle
filters represent the filtering distribution as a weighted sum of
Dirac-delta functions (\textbf{particles}):
\begin{equation}
  p(x) \approx \sum_{i=1}^{n} w^{(i)}\,\delta_{x^{(i)}}(x).
\end{equation}

\subsection{Importance Sampling (Section 3.8.1)}

To approximate a target $p(x)$ using samples from a \textbf{proposal}
$q(x)$:
\begin{equation}
  x^{(i)}\sim q, \qquad w^{(i)} = \frac{p(x^{(i)})}{q(x^{(i)})}.
\end{equation}
A large number of particles $n$ gives a better approximation.

\subsection{Resampling (Section 3.8.2)}

After dynamics and observation updates, some particles accumulate very
small weights (\textbf{particle degeneracy}). Resampling:
\begin{itemize}
  \item Takes weighted particles $\{w^{(i)},x^{(i)}\}_{i=1}^n$.
  \item Outputs new particles $\{x'^{(i)}\}$ with equal weights
    $w'^{(i)}=1/n$.
  \item Particles with high weights are duplicated; those with low
    weights are discarded.
  \item Algorithm: \textbf{low-variance resampling} (roulette wheel) in
    $O(n)$.
\end{itemize}

%============================================================
\section{Key Equations Reference}
%============================================================

\begin{center}
\renewcommand{\arraystretch}{1.7}
\begin{tabular}{ll}
\hline
\textbf{Filter} & \textbf{Key Update Equations} \\
\hline
KF predict & $\mu_{k+1|k}=A\mu_{k|k}+Bu_k$;\; $\Sigma_{k+1|k}=A\Sigma_{k|k}A^\top+R$ \\
KF update & $K=\Sigma_{k+1|k}C^\top(C\Sigma_{k+1|k}C^\top+Q)^{-1}$ \\
          & $\mu_{k+1|k+1}=\mu_{k+1|k}+K(y_{k+1}-C\mu_{k+1|k})$ \\
          & $\Sigma_{k+1|k+1}=(I-KC)\Sigma_{k+1|k}$ \\
EKF       & Replace $A\to\partial f/\partial x|_{\mu_{k|k}}$, $C\to\partial g/\partial x|_{\mu_{k+1|k}}$ \\
UKF       & Use sigma points; no Jacobians; 3rd-order accurate \\
PF        & Importance weights $w^{(i)}=p/q$; resample to avoid degeneracy \\
\hline
\end{tabular}
\end{center}

\end{document}
